\documentclass[preview]{standalone}

\usepackage{amsmath}
\usepackage{amssymb}
\usepackage{stellar}
\usepackage{definitions}

\begin{document}

\id{surfaces-and-flux-exercises}
\genpage

\section{Exercises}

\begin{snippetexercise}{surfaces-and-flux-exercises-ex1}{Helicoid}
    Let \(\Sigma\) be the surface parametrized by \[\sigma(u,v)=(u\cos v,u\sin v,v),\qquad u,v\in\realnumbers.\] Study the regularity of \(\Sigma\) and identify the coordinate curves \(u=\mathrm{const.}\) and \(v=\mathrm{const.}\). Determine whether \(\Sigma\) is equivalent to the graph surface of the function \(z=\arctan\left(\frac{y}{x}\right)\).
\end{snippetexercise}

\begin{snippetsolution}{surfaces-and-flux-exercises-ex1-sol}{Helicoid}
    We have
    \[
        \sigma_u(u,v)=(\cos v,\sin v,0),\qquad
        \sigma_v(u,v)=(-u\sin v,u\cos v,1).
    \]
    Hence
    \[
        \sigma_u(u,v)\times \sigma_v(u,v)
        =
        (\sin v,-\cos v,u),
    \]
    and therefore
    \[
        \|\sigma_u(u,v)\times \sigma_v(u,v)\|
        =
        \sqrt{1+u^2}>0
    \]
    for every \((u,v)\in\realnumbers^2\). Thus \(\sigma\) is a regular parametrization.

    Moreover, \(\sigma\) is injective. Indeed, if
    \[
        \sigma(u,v)=\sigma(\tilde u,\tilde v),
    \]
    then comparing the third coordinates gives \(v=\tilde v\). Then
    \[
        u\cos v=\tilde u\cos v,\qquad
        u\sin v=\tilde u\sin v,
    \]
    and since \((\cos v,\sin v)\neq (0,0)\), we get \(u=\tilde u\). Therefore \(\sigma\) is one-to-one. Its inverse on \(\Sigma\) is given by
    \[
        v=z,\qquad u=x\cos z+y\sin z,
    \]
    so \(\Sigma\) is a regular surface.

    If \(u=u_0\) is constant, then the corresponding coordinate curve is
    \[
        v\mapsto (u_0\cos v,u_0\sin v,v).
    \]
    For \(u_0\neq 0\), this is a helix around the \(z\)-axis with radius \(|u_0|\). For \(u_0=0\), it is the \(z\)-axis.

    If \(v=v_0\) is constant, then the corresponding coordinate curve is
    \[
        u\mapsto (u\cos v_0,u\sin v_0,v_0),
    \]
    which is a straight line contained in the horizontal plane \(z=v_0\). These lines are the rulings of the helicoid.

    Finally, \(\Sigma\) is not equivalent globally to the graph surface
    \[
        z=\arctan\left(\frac{y}{x}\right).
    \]
    Indeed, the helicoid is not the graph of a single-valued function over the \(xy\)-plane. For example, if \((x,y)\neq (0,0)\), the same point \((x,y)\) in the \(xy\)-plane corresponds to infinitely many points of the helicoid with different \(z\)-coordinates, since the angle can differ by integer multiples of \(\pi\), with the sign of \(u\) changing accordingly. On the other hand, the graph of
    \[
        z=\arctan\left(\frac{y}{x}\right)
    \]
    assigns only one value of \(z\) to each point of its domain. Thus it can describe only a local branch of the helicoid, not the whole surface \(\Sigma\).
\end{snippetsolution}

\begin{snippetexercise}{surfaces-and-flux-exercises-ex2}{Catenoid}
    Write a parametrization of the catenoid, the surface described by rotating the curve \(x=\cosh z\) in the \(xz\)-plane around the \(z\)-axis. Verify that the catenoid is a regular surface and determine its unit normal vector and area element.
\end{snippetexercise}

\begin{snippetsolution}{surfaces-and-flux-exercises-ex2-sol}{Catenoid}
    Rotating the curve \(x=\cosh z\) around the \(z\)-axis gives the parametrization
    \[
        \sigma(u,v)=(\cosh u\cos v,\cosh u\sin v,u),
        \qquad u\in\realnumbers,\ v\in\realnumbers.
    \]
    Equivalently, one may take \(v\in(0,2\pi)\) to obtain a local chart.

    We compute
    \[
        \sigma_u(u,v)=(\sinh u\cos v,\sinh u\sin v,1),
    \]
    and
    \[
        \sigma_v(u,v)=(-\cosh u\sin v,\cosh u\cos v,0).
    \]
    Hence
    \[
        \sigma_u(u,v)\times\sigma_v(u,v)
        =
        (-\cosh u\cos v,-\cosh u\sin v,\sinh u\cosh u).
    \]
    Therefore
    \[
        \|\sigma_u(u,v)\times\sigma_v(u,v)\|
        =
        \cosh u\sqrt{1+\sinh^2u}
        =
        \cosh^2 u>0
    \]
    for every \(u\in\realnumbers\). Thus the parametrization is regular. Since the catenoid can be covered by such regular local charts, it is a regular surface.

    A unit normal vector is given by
    \[
        N(u,v)
        =
        \frac{\sigma_u(u,v)\times\sigma_v(u,v)}
        {\|\sigma_u(u,v)\times\sigma_v(u,v)\|}
        =
        \left(
            -\frac{\cos v}{\cosh u},
            -\frac{\sin v}{\cosh u},
            \tanh u
        \right).
    \]
    The opposite orientation is obtained by taking \(-N(u,v)\).

    Finally, the area element is
    \[
        dS
        =
        \|\sigma_u(u,v)\times\sigma_v(u,v)\|\,du\,dv
        =
        \cosh^2 u\,du\,dv.
    \]
\end{snippetsolution}

\begin{snippetexercise}{surfaces-and-flux-exercises-ex3}{Pappus theorem}
    Prove Pappus's theorem: let \(\Gamma\) be a regular arc contained in the \(xz\)-plane and not intersecting the \(z\)-axis. Let \(\Sigma\) be the surface of revolution obtained by a full rotation of \(\Gamma\) around the \(z\)-axis. Then \[A(\Sigma)=2\pi \bar{x}L(\Gamma),\] where \(L(\Gamma)\) is the length of \(\Gamma\) and \(\bar{x}\) is the \(x\)-coordinate of the centroid of \(\Gamma\).
\end{snippetexercise}

\begin{snippetsolution}{surfaces-and-flux-exercises-ex3-sol}{Pappus theorem}
    Let
    \[
        \gamma(t)=(x(t),0,z(t)),\qquad t\in[a,b],
    \]
    be a regular parametrization of \(\Gamma\). Since \(\Gamma\) does not intersect the \(z\)-axis, we have \(x(t)\neq 0\) for every \(t\in[a,b]\). Since \([a,b]\) is connected, \(x(t)\) has constant sign. Up to replacing \(x\) by \(-x\), we may assume \(x(t)>0\).

    The surface of revolution obtained by rotating \(\Gamma\) around the \(z\)-axis is parametrized by
    \[
        \sigma(t,\theta)
        =
        (x(t)\cos\theta,x(t)\sin\theta,z(t)),
        \qquad
        t\in[a,b],\ \theta\in[0,2\pi].
    \]
    We compute
    \[
        \sigma_t(t,\theta)
        =
        (x'(t)\cos\theta,x'(t)\sin\theta,z'(t)),
    \]
    and
    \[
        \sigma_\theta(t,\theta)
        =
        (-x(t)\sin\theta,x(t)\cos\theta,0).
    \]
    Hence
    \[
        \sigma_t(t,\theta)\times\sigma_\theta(t,\theta)
        =
        \bigl(
            -x(t)z'(t)\cos\theta,
            -x(t)z'(t)\sin\theta,
            x(t)x'(t)
        \bigr).
    \]
    Therefore
    \[
        \|\sigma_t(t,\theta)\times\sigma_\theta(t,\theta)\|
        =
        x(t)\sqrt{(x'(t))^2+(z'(t))^2}.
    \]
    In particular, since \(\gamma\) is regular and \(x(t)>0\), the parametrization is regular.

    The area of \(\Sigma\) is then
    \[
        A(\Sigma)
        =
        \int_0^{2\pi}\int_a^b
        x(t)\sqrt{(x'(t))^2+(z'(t))^2}
        \,dt\,d\theta.
    \]
    Thus
    \[
        A(\Sigma)
        =
        2\pi
        \int_a^b
        x(t)\sqrt{(x'(t))^2+(z'(t))^2}
        \,dt.
    \]
    Since the arc-length element of \(\Gamma\) is
    \[
        ds=\sqrt{(x'(t))^2+(z'(t))^2}\,dt,
    \]
    we get
    \[
        A(\Sigma)
        =
        2\pi\int_\Gamma x\,ds.
    \]

    By definition, the \(x\)-coordinate of the centroid of \(\Gamma\) is
    \[
        \bar{x}
        =
        \frac{1}{L(\Gamma)}\int_\Gamma x\,ds.
    \]
    Therefore
    \[
        \int_\Gamma x\,ds=\bar{x}L(\Gamma),
    \]
    and so
    \[
        A(\Sigma)=2\pi \bar{x}L(\Gamma).
    \]
    If the arc lies in the half-plane \(x<0\), the same argument gives
    \[
        A(\Sigma)=2\pi |\bar{x}|L(\Gamma),
    \]
    namely the distance of the centroid from the axis times the length of the arc.
\end{snippetsolution}

\begin{snippetexercise}{surfaces-and-flux-exercises-ex4}{Cone inside a cylinder}
    Compute the area of the portion of the cone with equation \(x^2=y^2+z^2\) inside the cylinder with equation \(x^2+y^2=1\).
\end{snippetexercise}

\begin{snippetsolution}{surfaces-and-flux-exercises-ex4-sol}{Cone inside a cylinder}
    The cone
    \[
        x^2=y^2+z^2
    \]
    has two nappes. We first compute the area of the nappe \(x\geq 0\), and then double the result.

    A parametrization of the nappe \(x\geq 0\) is
    \[
        \sigma(r,\theta)=(r,r\cos\theta,r\sin\theta),
        \qquad r\geq 0,\ \theta\in[0,2\pi].
    \]
    On the cone we have
    \[
        x^2+y^2=r^2+r^2\cos^2\theta
        =
        r^2(1+\cos^2\theta).
    \]
    The condition of being inside the cylinder \(x^2+y^2=1\) is therefore
    \[
        r^2(1+\cos^2\theta)\leq 1,
    \]
    that is
    \[
        0\leq r\leq \frac{1}{\sqrt{1+\cos^2\theta}}.
    \]

    We compute
    \[
        \sigma_r(r,\theta)=(1,\cos\theta,\sin\theta),
    \]
    and
    \[
        \sigma_\theta(r,\theta)=(0,-r\sin\theta,r\cos\theta).
    \]
    Hence
    \[
        \sigma_r(r,\theta)\times\sigma_\theta(r,\theta)
        =
        (r,-r\cos\theta,-r\sin\theta),
    \]
    and so
    \[
        \|\sigma_r(r,\theta)\times\sigma_\theta(r,\theta)\|
        =
        \sqrt{2}\,r.
    \]

    Thus the area of one nappe inside the cylinder is
    \[
        A_+
        =
        \int_0^{2\pi}
        \int_0^{\frac{1}{\sqrt{1+\cos^2\theta}}}
        \sqrt{2}\,r\,dr\,d\theta.
    \]
    Therefore
    \[
        A_+
        =
        \frac{\sqrt{2}}{2}
        \int_0^{2\pi}
        \frac{1}{1+\cos^2\theta}
        \,d\theta.
    \]
    Using the standard integral
    \[
        \int_0^{2\pi}\frac{d\theta}{1+\cos^2\theta}
        =
        \sqrt{2}\pi,
    \]
    we obtain
    \[
        A_+=\pi.
    \]

    Since the cone has two symmetric nappes, the total area is
    \[
        A=2A_+=2\pi.
    \]
\end{snippetsolution}

\begin{snippetexercise}{surfaces-and-flux-exercises-ex5}{Cylinder inside a cylinder}
    Let \(a>0\). Compute the area of the portion of the cylinder with equation \(y^2+z^2=a^2\) inside the cylinder with equation \(x^2+y^2=a^2\).
\end{snippetexercise}

\begin{snippetsolution}{surfaces-and-flux-exercises-ex5-sol}{Cylinder inside a cylinder}
    We parametrize the cylinder
    \[
        y^2+z^2=a^2
    \]
    by
    \[
        \sigma(x,\theta)=(x,a\cos\theta,a\sin\theta),
        \qquad x\in\realnumbers,\ \theta\in[0,2\pi].
    \]
    Then
    \[
        \sigma_x(x,\theta)=(1,0,0),
    \]
    and
    \[
        \sigma_\theta(x,\theta)=(0,-a\sin\theta,a\cos\theta).
    \]
    Hence
    \[
        \sigma_x(x,\theta)\times\sigma_\theta(x,\theta)
        =
        (0,-a\cos\theta,-a\sin\theta),
    \]
    so
    \[
        \|\sigma_x(x,\theta)\times\sigma_\theta(x,\theta)\|=a.
    \]

    We now impose the condition that the point lies inside the cylinder
    \[
        x^2+y^2=a^2.
    \]
    Since on the first cylinder we have \(y=a\cos\theta\), this gives
    \[
        x^2+a^2\cos^2\theta\leq a^2.
    \]
    Therefore
    \[
        x^2\leq a^2\sin^2\theta,
    \]
    and so
    \[
        -a|\sin\theta|\leq x\leq a|\sin\theta|.
    \]

    Thus the required area is
    \[
        A
        =
        \int_0^{2\pi}
        \int_{-a|\sin\theta|}^{a|\sin\theta|}
        a\,dx\,d\theta.
    \]
    Hence
    \[
        A
        =
        2a^2\int_0^{2\pi}|\sin\theta|\,d\theta.
    \]
    Since
    \[
        \int_0^{2\pi}|\sin\theta|\,d\theta=4,
    \]
    we obtain
    \[
        A=8a^2.
    \]
\end{snippetsolution}

\begin{snippetexercise}{surfaces-and-flux-exercises-ex6}{Centroid of a spherical portion}
    Compute the coordinates of the centroid of the portion of the sphere with equation \(x^2+y^2+z^2=R^2\), \(x\geq 0\), \(y\geq 0\).
\end{snippetexercise}

\begin{snippetsolution}{surfaces-and-flux-exercises-ex6-sol}{Centroid of a spherical portion}
    We consider the spherical coordinates
    \[
        \sigma(\varphi,\theta)
        =
        (R\sin\varphi\cos\theta,
        R\sin\varphi\sin\theta,
        R\cos\varphi),
    \]
    with
    \[
        0\leq \varphi\leq \pi,
        \qquad
        0\leq \theta\leq \frac{\pi}{2}.
    \]
    These bounds describe the portion of the sphere satisfying \(x\geq 0\) and \(y\geq 0\).

    The area element on the sphere is
    \[
        dS=R^2\sin\varphi\,d\varphi\,d\theta.
    \]
    Hence the area of the surface is
    \[
        A
        =
        \int_0^{\frac{\pi}{2}}\int_0^\pi
        R^2\sin\varphi\,d\varphi\,d\theta
        =
        \pi R^2.
    \]

    By symmetry with respect to the plane \(x=y\), we have
    \[
        \bar{x}=\bar{y}.
    \]
    Moreover, by symmetry with respect to the plane \(z=0\), we have
    \[
        \bar{z}=0.
    \]

    Now
    \[
        \bar{x}
        =
        \frac{1}{A}\int_\Sigma x\,dS.
    \]
    Therefore
    \[
        \int_\Sigma x\,dS
        =
        \int_0^{\frac{\pi}{2}}\int_0^\pi
        R\sin\varphi\cos\theta
        R^2\sin\varphi
        \,d\varphi\,d\theta.
    \]
    Thus
    \[
        \int_\Sigma x\,dS
        =
        R^3
        \left(\int_0^{\frac{\pi}{2}}\cos\theta\,d\theta\right)
        \left(\int_0^\pi \sin^2\varphi\,d\varphi\right)
        =
        R^3\cdot 1\cdot \frac{\pi}{2}
        =
        \frac{\pi R^3}{2}.
    \]
    Hence
    \[
        \bar{x}
        =
        \frac{1}{\pi R^2}\frac{\pi R^3}{2}
        =
        \frac{R}{2}.
    \]
    Similarly,
    \[
        \bar{y}=\frac{R}{2}.
    \]

    Therefore the centroid is
    \[
        \left(\frac{R}{2},\frac{R}{2},0\right).
    \]
\end{snippetsolution}

\begin{snippetexercise}{surfaces-and-flux-exercises-ex7}{Torricelli trumpet}
    Parametrize the surface \(\Sigma\) described by rotating the curve \(y=1/x\), for \(x\in[1,+\infty)\), around the \(x\)-axis. Verify that \(\Sigma\) is a regular surface and determine its area element. Verify that, if \(\Sigma\) were a real object, an infinite amount of paint would be needed to paint its exterior face, but a finite amount of paint would be enough to fill it. Let \(\Sigma_0\) be the portion of \(\Sigma\) with \(x\in[1,2]\). Compute the integral \[\iint_{\Sigma_0}\frac{1}{x^4}\,\mathrm{d}S.\]
\end{snippetexercise}

\begin{snippetsolution}{surfaces-and-flux-exercises-ex7-sol}{Torricelli trumpet}
    Rotating the curve \(y=1/x\), with \(x\in[1,+\infty)\), around the \(x\)-axis gives the parametrization
    \[
        \sigma(u,v)
        =
        \left(u,\frac{\cos v}{u},\frac{\sin v}{u}\right),
        \qquad
        u\in[1,+\infty),\ v\in[0,2\pi].
    \]
    We compute
    \[
        \sigma_u(u,v)
        =
        \left(1,-\frac{\cos v}{u^2},-\frac{\sin v}{u^2}\right),
    \]
    and
    \[
        \sigma_v(u,v)
        =
        \left(0,-\frac{\sin v}{u},\frac{\cos v}{u}\right).
    \]
    Hence
    \[
        \sigma_u(u,v)\times\sigma_v(u,v)
        =
        \left(
            -\frac{1}{u^3},
            -\frac{\cos v}{u},
            -\frac{\sin v}{u}
        \right),
    \]
    and therefore
    \[
        \|\sigma_u(u,v)\times\sigma_v(u,v)\|
        =
        \sqrt{\frac{1}{u^6}+\frac{1}{u^2}}
        =
        \frac{1}{u}\sqrt{1+\frac{1}{u^4}}>0.
    \]
    Thus \(\Sigma\) is a regular surface. Its area element is
    \[
        dS
        =
        \frac{1}{u}\sqrt{1+\frac{1}{u^4}}\,du\,dv.
    \]

    The area of \(\Sigma\) is
    \[
        A(\Sigma)
        =
        \int_0^{2\pi}\int_1^{+\infty}
        \frac{1}{u}\sqrt{1+\frac{1}{u^4}}
        \,du\,dv.
    \]
    Since
    \[
        \frac{1}{u}\sqrt{1+\frac{1}{u^4}}
        \geq
        \frac{1}{u},
    \]
    we have
    \[
        A(\Sigma)
        \geq
        2\pi\int_1^{+\infty}\frac{1}{u}\,du
        =
        +\infty.
    \]
    Therefore an infinite amount of paint would be needed to paint the exterior face.

    On the other hand, the volume enclosed by the trumpet is
    \[
        V
        =
        \pi\int_1^{+\infty}\frac{1}{u^2}\,du
        =
        \pi.
    \]
    Hence a finite amount of paint would be enough to fill it.

    Now let \(\Sigma_0\) be the portion with \(1\leq x\leq 2\), that is \(1\leq u\leq 2\). Since \(x=u\), we get
    \[
        \iint_{\Sigma_0}\frac{1}{x^4}\,dS
        =
        \int_0^{2\pi}\int_1^2
        \frac{1}{u^4}
        \frac{1}{u}\sqrt{1+\frac{1}{u^4}}
        \,du\,dv.
    \]
    Therefore
    \[
        \iint_{\Sigma_0}\frac{1}{x^4}\,dS
        =
        2\pi\int_1^2
        \frac{1}{u^5}\sqrt{1+\frac{1}{u^4}}
        \,du.
    \]
    With the substitution
    \[
        t=1+\frac{1}{u^4},
        \qquad
        dt=-\frac{4}{u^5}\,du,
    \]
    we obtain
    \[
        \iint_{\Sigma_0}\frac{1}{x^4}\,dS
        =
        \frac{\pi}{2}\int_{\frac{17}{16}}^2 t^{\frac12}\,dt.
    \]
    Thus
    \[
        \iint_{\Sigma_0}\frac{1}{x^4}\,dS
        =
        \frac{\pi}{3}
        \left[
            t^{\frac32}
        \right]_{\frac{17}{16}}^2
        =
        \frac{\pi}{3}
        \left(
            2\sqrt{2}-\frac{17\sqrt{17}}{64}
        \right).
    \]
\end{snippetsolution}

\begin{snippetexercise}{surfaces-and-flux-exercises-ex8}{Flux through a parametrized surface}
    Compute the flux of the field \[\mathbf{F}(x,y,z)=\mathbf{i}+\mathbf{k}\] through the surface \(\Sigma\) parametrized by \[\sigma(u,v)=\left(u^2,\sqrt{2}uv,v^2\right),\qquad (u,v)\in T,\] where \[T=\{(u,v):1<u^2+v^2<2,\ u<v\}.\]
\end{snippetexercise}

\begin{snippetsolution}{surfaces-and-flux-exercises-ex8-sol}{Flux through a parametrized surface}
    The orientation is the one induced by the parametrization. We have
    \[
        \sigma_u(u,v)
        =
        (2u,\sqrt{2}v,0),
    \]
    and
    \[
        \sigma_v(u,v)
        =
        (0,\sqrt{2}u,2v).
    \]
    Hence
    \[
        \sigma_u(u,v)\times\sigma_v(u,v)
        =
        (2\sqrt{2}v^2,-4uv,2\sqrt{2}u^2).
    \]
    Since \((u,v)\in T\) implies \(u^2+v^2>1\), this vector never vanishes on \(T\). Thus the parametrization is regular on \(T\).

    Moreover,
    \[
        \mathbf{F}(\sigma(u,v))
        =
        (1,0,1).
    \]
    Therefore
    \[
        \mathbf{F}(\sigma(u,v))\cdot
        \left(\sigma_u(u,v)\times\sigma_v(u,v)\right)
        =
        2\sqrt{2}v^2+2\sqrt{2}u^2
        =
        2\sqrt{2}(u^2+v^2).
    \]
    Hence the flux is
    \[
        \Phi
        =
        \iint_T 2\sqrt{2}(u^2+v^2)\,du\,dv.
    \]

    We use polar coordinates
    \[
        u=r\cos\theta,\qquad v=r\sin\theta.
    \]
    The condition
    \[
        1<u^2+v^2<2
    \]
    becomes
    \[
        1<r^2<2,
    \]
    so
    \[
        1<r<\sqrt{2}.
    \]
    The condition \(u<v\) becomes
    \[
        \cos\theta<\sin\theta,
    \]
    which corresponds to
    \[
        \frac{\pi}{4}<\theta<\frac{5\pi}{4}.
    \]
    Therefore
    \[
        \Phi
        =
        2\sqrt{2}
        \int_{\frac{\pi}{4}}^{\frac{5\pi}{4}}
        \int_1^{\sqrt{2}}
        r^2r\,dr\,d\theta.
    \]
    Thus
    \[
        \Phi
        =
        2\sqrt{2}
        \left(\int_{\frac{\pi}{4}}^{\frac{5\pi}{4}}d\theta\right)
        \left(\int_1^{\sqrt{2}}r^3\,dr\right).
    \]
    Since
    \[
        \int_{\frac{\pi}{4}}^{\frac{5\pi}{4}}d\theta=\pi
    \]
    and
    \[
        \int_1^{\sqrt{2}}r^3\,dr
        =
        \frac{1}{4}\left((\sqrt{2})^4-1\right)
        =
        \frac{3}{4},
    \]
    we obtain
    \[
        \Phi
        =
        2\sqrt{2}\pi\frac{3}{4}
        =
        \frac{3\sqrt{2}\pi}{2}.
    \]
    With the opposite orientation, the flux would be
    \[
        -\frac{3\sqrt{2}\pi}{2}.
    \]
\end{snippetsolution}

\begin{snippetexercise}{surfaces-and-flux-exercises-ex9}{Flux through a graph}
    Compute the flux of the field \[\mathbf{F}(x,y,z)=y\mathbf{i}+x\mathbf{j}\] through the surface \(\Sigma\) with equation \(z=xy\), which projects onto \[T=\left\{(x,y):1\leq x\leq\sqrt{2-y^2}\right\}.\]
\end{snippetexercise}

\begin{snippetsolution}{surfaces-and-flux-exercises-ex9-sol}{Flux through a graph}
    We parametrize the graph \(z=xy\) by
    \[
        \sigma(x,y)=(x,y,xy),
        \qquad (x,y)\in T.
    \]
    Then
    \[
        \sigma_x(x,y)=(1,0,y),
        \qquad
        \sigma_y(x,y)=(0,1,x),
    \]
    and therefore
    \[
        \sigma_x(x,y)\times\sigma_y(x,y)=(-y,-x,1).
    \]
    This corresponds to the upward orientation.

    Since
    \[
        \mathbf{F}(\sigma(x,y))=(y,x,0),
    \]
    we get
    \[
        \mathbf{F}(\sigma(x,y))\cdot
        \left(\sigma_x(x,y)\times\sigma_y(x,y)\right)
        =
        -y^2-x^2.
    \]
    Hence the flux is
    \[
        \Phi
        =
        -\iint_T (x^2+y^2)\,dx\,dy.
    \]

    The region \(T\) is given by
    \[
        T=\{(x,y):x^2+y^2\leq 2,\ x\geq 1\}.
    \]
    In polar coordinates,
    \[
        x=r\cos\theta,\qquad y=r\sin\theta,
    \]
    the condition \(x\geq 1\) becomes
    \[
        r\cos\theta\geq 1.
    \]
    Therefore
    \[
        -\frac{\pi}{4}\leq \theta\leq \frac{\pi}{4},
        \qquad
        \frac{1}{\cos\theta}\leq r\leq \sqrt{2}.
    \]
    Thus
    \[
        \Phi
        =
        -\int_{-\frac{\pi}{4}}^{\frac{\pi}{4}}
        \int_{\frac{1}{\cos\theta}}^{\sqrt{2}}
        r^2 r\,dr\,d\theta.
    \]
    Hence
    \[
        \Phi
        =
        -\frac{1}{4}
        \int_{-\frac{\pi}{4}}^{\frac{\pi}{4}}
        \left(4-\sec^4\theta\right)
        \,d\theta.
    \]
    Since
    \[
        \int \sec^4\theta\,d\theta
        =
        \tan\theta+\frac{\tan^3\theta}{3},
    \]
    we have
    \[
        \int_{-\frac{\pi}{4}}^{\frac{\pi}{4}}\sec^4\theta\,d\theta
        =
        \frac{8}{3}.
    \]
    Therefore
    \[
        \Phi
        =
        -\frac{1}{4}
        \left(2\pi-\frac{8}{3}\right)
        =
        \frac{2}{3}-\frac{\pi}{2}.
    \]
    With the opposite orientation, the flux would be
    \[
        \frac{\pi}{2}-\frac{2}{3}.
    \]
\end{snippetsolution}

\begin{snippetexercise}{surfaces-and-flux-exercises-ex10}{Flux through a cylindrical surface}
    Compute the flux of the field \[\mathbf{F}(x,y,z)=ye^{-x}\mathbf{j}\] through the surface \[\Sigma=\{(x,y,z):x^2+y^2=1,\ 0\leq z\leq y\}.\] Choose the outward normal to the cylindrical surface.
\end{snippetexercise}

\begin{snippetsolution}{surfaces-and-flux-exercises-ex10-sol}{Flux through a cylindrical surface}
    Since
    \[
        x^2+y^2=1
    \]
    and
    \[
        0\leq z\leq y,
    \]
    we must have \(y\geq 0\). Thus we parametrize the relevant portion of the cylinder by
    \[
        \sigma(\theta,z)=(\cos\theta,\sin\theta,z),
        \qquad
        0\leq \theta\leq \pi,\quad
        0\leq z\leq \sin\theta.
    \]
    We compute
    \[
        \sigma_\theta(\theta,z)=(-\sin\theta,\cos\theta,0),
        \qquad
        \sigma_z(\theta,z)=(0,0,1).
    \]
    Hence
    \[
        \sigma_\theta(\theta,z)\times\sigma_z(\theta,z)
        =
        (\cos\theta,\sin\theta,0),
    \]
    which is the outward normal vector to the cylinder.

    The vector field is
    \[
        \mathbf{F}(x,y,z)=(0,ye^{-x},0).
    \]
    Therefore
    \[
        \mathbf{F}(\sigma(\theta,z))
        =
        (0,\sin\theta e^{-\cos\theta},0).
    \]
    Hence
    \[
        \mathbf{F}(\sigma(\theta,z))\cdot
        \left(\sigma_\theta(\theta,z)\times\sigma_z(\theta,z)\right)
        =
        \sin^2\theta e^{-\cos\theta}.
    \]
    The flux is therefore
    \[
        \Phi
        =
        \int_0^\pi
        \int_0^{\sin\theta}
        \sin^2\theta e^{-\cos\theta}
        \,dz\,d\theta.
    \]
    Thus
    \[
        \Phi
        =
        \int_0^\pi
        \sin^3\theta e^{-\cos\theta}
        \,d\theta.
    \]
    With the substitution
    \[
        t=\cos\theta,
        \qquad
        dt=-\sin\theta\,d\theta,
    \]
    we obtain
    \[
        \Phi
        =
        \int_{-1}^1
        (1-t^2)e^{-t}\,dt.
    \]
    Since
    \[
        \int (1-t^2)e^{-t}\,dt
        =
        e^{-t}(t+1)^2,
    \]
    it follows that
    \[
        \Phi
        =
        \left[e^{-t}(t+1)^2\right]_{-1}^1
        =
        \frac{4}{e}.
    \]
\end{snippetsolution}

\begin{snippetexercise}{surfaces-and-flux-exercises-ex11}{Green-Gauss theorem on a square}
    Compute the integral \[\int_{\gamma}\left(y^2\,\mathrm{d}x+x\,\mathrm{d}y\right),\] where \(\gamma\) is the boundary of the square with vertices \((0,0)\), \((1,0)\), \((1,1)\), \((0,1)\), traversed counterclockwise. Verify the result by applying the Green-Gauss theorem.
\end{snippetexercise}

\begin{snippetsolution}{surfaces-and-flux-exercises-ex11-sol}{Green-Gauss theorem on a square}
    Let
    \[
        P(x,y)=y^2,
        \qquad
        Q(x,y)=x.
    \]
    We first compute the integral directly, splitting \(\gamma\) into the four sides of the square.

    On the side from \((0,0)\) to \((1,0)\), we have \(y=0\), \(dy=0\), and therefore
    \[
        \int y^2\,dx+x\,dy=0.
    \]
    On the side from \((1,0)\) to \((1,1)\), we have \(x=1\), \(dx=0\), and hence
    \[
        \int y^2\,dx+x\,dy
        =
        \int_0^1 dy
        =
        1.
    \]
    On the side from \((1,1)\) to \((0,1)\), we have \(y=1\), \(dy=0\), and \(x\) goes from \(1\) to \(0\). Hence
    \[
        \int y^2\,dx+x\,dy
        =
        \int_1^0 dx
        =
        -1.
    \]
    On the side from \((0,1)\) to \((0,0)\), we have \(x=0\), \(dx=0\), and therefore
    \[
        \int y^2\,dx+x\,dy=0.
    \]
    Summing the four contributions, we get
    \[
        \int_{\gamma}\left(y^2\,dx+x\,dy\right)
        =
        0+1-1+0
        =
        0.
    \]

    We now verify this using the Green-Gauss theorem. Since \(\gamma\) is positively oriented,
    \[
        \int_{\gamma}P\,dx+Q\,dy
        =
        \iint_{[0,1]\times[0,1]}
        \left(
            \frac{\partial Q}{\partial x}
            -
            \frac{\partial P}{\partial y}
        \right)
        \,dx\,dy.
    \]
    Since
    \[
        \frac{\partial Q}{\partial x}=1,
        \qquad
        \frac{\partial P}{\partial y}=2y,
    \]
    we obtain
    \[
        \int_{\gamma}\left(y^2\,dx+x\,dy\right)
        =
        \int_0^1\int_0^1 (1-2y)\,dx\,dy.
    \]
    Therefore
    \[
        \int_{\gamma}\left(y^2\,dx+x\,dy\right)
        =
        \int_0^1(1-2y)\,dy
        =
        \left[y-y^2\right]_0^1
        =
        0.
    \]
\end{snippetsolution}

\begin{snippetexercise}{surfaces-and-flux-exercises-ex12}{Double integral as a line integral}
    Compute \[\iint_D x^2\,\mathrm{d}x\mathrm{d}y,\] where \[D=\{(x,y):1\leq x^2+y^2\leq 2\},\] by transforming the double integral into a line integral. Verify the result by computing the integral directly.
\end{snippetexercise}

\begin{snippetsolution}{surfaces-and-flux-exercises-ex12-sol}{Double integral as a line integral}
    We want to write
    \[
        \iint_D x^2\,dx\,dy
    \]
    as a line integral. By Green-Gauss theorem, it is enough to find \(P,Q\) such that
    \[
        \frac{\partial Q}{\partial x}
        -
        \frac{\partial P}{\partial y}
        =
        x^2.
    \]
    We can choose
    \[
        P(x,y)=0,
        \qquad
        Q(x,y)=\frac{x^3}{3}.
    \]
    Hence
    \[
        \iint_D x^2\,dx\,dy
        =
        \int_{\partial D}\frac{x^3}{3}\,dy,
    \]
    where \(\partial D\) is positively oriented.

    The boundary of \(D\) consists of two circles:
    \[
        C_2:x^2+y^2=2,
        \qquad
        C_1:x^2+y^2=1.
    \]
    The outer circle \(C_2\) is traversed counterclockwise, while the inner circle \(C_1\) is traversed clockwise.

    For a circle of radius \(r\), parametrized counterclockwise by
    \[
        \gamma(t)=(r\cos t,r\sin t),
        \qquad 0\leq t\leq 2\pi,
    \]
    we have
    \[
        dy=r\cos t\,dt.
    \]
    Therefore
    \[
        \int_\gamma \frac{x^3}{3}\,dy
        =
        \int_0^{2\pi}
        \frac{r^3\cos^3 t}{3}r\cos t\,dt
        =
        \frac{r^4}{3}\int_0^{2\pi}\cos^4 t\,dt.
    \]
    Since
    \[
        \int_0^{2\pi}\cos^4 t\,dt=\frac{3\pi}{4},
    \]
    we get
    \[
        \int_\gamma \frac{x^3}{3}\,dy
        =
        \frac{\pi r^4}{4}.
    \]

    Thus the contribution of the outer boundary is
    \[
        \frac{\pi(\sqrt{2})^4}{4}=\pi,
    \]
    while the contribution of the inner boundary, which has the opposite orientation, is
    \[
        -\frac{\pi}{4}.
    \]
    Hence
    \[
        \iint_D x^2\,dx\,dy
        =
        \pi-\frac{\pi}{4}
        =
        \frac{3\pi}{4}.
    \]

    We now verify the result directly using polar coordinates. Since
    \[
        x=r\cos\theta,
        \qquad
        dx\,dy=r\,dr\,d\theta,
    \]
    and
    \[
        1\leq r\leq \sqrt{2},
        \qquad
        0\leq \theta\leq 2\pi,
    \]
    we obtain
    \[
        \iint_D x^2\,dx\,dy
        =
        \int_0^{2\pi}\int_1^{\sqrt{2}}
        r^2\cos^2\theta\,r\,dr\,d\theta.
    \]
    Therefore
    \[
        \iint_D x^2\,dx\,dy
        =
        \left(\int_0^{2\pi}\cos^2\theta\,d\theta\right)
        \left(\int_1^{\sqrt{2}}r^3\,dr\right).
    \]
    Since
    \[
        \int_0^{2\pi}\cos^2\theta\,d\theta=\pi
    \]
    and
    \[
        \int_1^{\sqrt{2}}r^3\,dr
        =
        \frac{1}{4}\left((\sqrt{2})^4-1\right)
        =
        \frac{3}{4},
    \]
    we get
    \[
        \iint_D x^2\,dx\,dy
        =
        \frac{3\pi}{4}.
    \]
\end{snippetsolution}

\begin{snippetexercise}{surfaces-and-flux-exercises-ex13}{Area formula for a plane curve}
    Let \(\Gamma\) be a closed, simple, regular, positively oriented plane curve parametrized by \[\gamma(t)=(x(t),y(t)),\qquad t\in[a,b].\] If \(\Gamma\) is the boundary of the domain \(D\), prove that \[|D|=\frac{1}{2}\int_a^b
        \begin{vmatrix} x(t) & y(t) \\ x'(t) & y'(t)
    \end{vmatrix}\,\mathrm{d}t.\]
\end{snippetexercise}

\begin{snippetsolution}{surfaces-and-flux-exercises-ex13-sol}{Area formula for a plane curve}
    By the Green-Gauss theorem, if \(P,Q\in C^1\) and \(\Gamma=\partial D\) is positively oriented, then
    \[
        \int_\Gamma P\,dx+Q\,dy
        =
        \iint_D
        \left(
            \frac{\partial Q}{\partial x}
            -
            \frac{\partial P}{\partial y}
        \right)
        \,dx\,dy.
    \]
    We choose
    \[
        P(x,y)=-\frac{y}{2},
        \qquad
        Q(x,y)=\frac{x}{2}.
    \]
    Then
    \[
        \frac{\partial Q}{\partial x}
        -
        \frac{\partial P}{\partial y}
        =
        \frac{1}{2}-\left(-\frac{1}{2}\right)
        =
        1.
    \]
    Hence
    \[
        |D|
        =
        \iint_D 1\,dx\,dy
        =
        \int_\Gamma
        \left(
            -\frac{y}{2}\,dx+\frac{x}{2}\,dy
        \right).
    \]
    Since \(\Gamma\) is parametrized by
    \[
        \gamma(t)=(x(t),y(t)),
        \qquad t\in[a,b],
    \]
    we have
    \[
        dx=x'(t)\,dt,
        \qquad
        dy=y'(t)\,dt.
    \]
    Therefore
    \[
        |D|
        =
        \frac{1}{2}
        \int_a^b
        \left(
            x(t)y'(t)-y(t)x'(t)
        \right)
        \,dt.
    \]
    Since
    \[
        x(t)y'(t)-y(t)x'(t)
        =
        \begin{vmatrix}
            x(t) & y(t) \\
            x'(t) & y'(t)
        \end{vmatrix},
    \]
    we obtain
    \[
        |D|
        =
        \frac{1}{2}\int_a^b
        \begin{vmatrix}
            x(t) & y(t) \\
            x'(t) & y'(t)
        \end{vmatrix}
        \,dt.
    \]
\end{snippetsolution}

\begin{snippetexercise}{surfaces-and-flux-exercises-ex14}{Area enclosed by an astroid}
    Compute the area enclosed by the astroid \[\gamma(t)=\left(\cos^3 t,\sin^3 t\right),\qquad t\in[0,2\pi].\]
\end{snippetexercise}

\begin{snippetsolution}{surfaces-and-flux-exercises-ex14-sol}{Area enclosed by an astroid}
    We use the area formula for a positively oriented closed plane curve:
    \[
        |D|
        =
        \frac{1}{2}\int_0^{2\pi}
        \begin{vmatrix}
            x(t) & y(t) \\
            x'(t) & y'(t)
        \end{vmatrix}
        \,dt.
    \]
    Here
    \[
        x(t)=\cos^3 t,
        \qquad
        y(t)=\sin^3 t.
    \]
    Thus
    \[
        x'(t)=-3\cos^2 t\sin t,
        \qquad
        y'(t)=3\sin^2 t\cos t.
    \]
    Therefore
    \[
        x(t)y'(t)-y(t)x'(t)
        =
        3\cos^4 t\sin^2 t
        +
        3\sin^4 t\cos^2 t.
    \]
    Hence
    \[
        x(t)y'(t)-y(t)x'(t)
        =
        3\sin^2 t\cos^2 t
        =
        \frac{3}{4}\sin^2(2t).
    \]
    It follows that
    \[
        |D|
        =
        \frac{1}{2}
        \int_0^{2\pi}
        3\sin^2 t\cos^2 t
        \,dt.
    \]
    Since
    \[
        \sin^2 t\cos^2 t=\frac{1}{8}\left(1-\cos(4t)\right),
    \]
    we obtain
    \[
        |D|
        =
        \frac{3}{2}
        \int_0^{2\pi}
        \frac{1}{8}\left(1-\cos(4t)\right)
        \,dt.
    \]
    Therefore
    \[
        |D|
        =
        \frac{3}{16}\cdot 2\pi
        =
        \frac{3\pi}{8}.
    \]
\end{snippetsolution}

\begin{snippetexercise}{surfaces-and-flux-exercises-ex15}{Stokes theorem on a circle}
    Compute the line integral of the field \[\mathbf{F}(x,y,z)=\left(x^2,x,y\right)\] along the circle in the plane \(z=0\) with equation \(x^2+y^2=4\), traversed counterclockwise. Verify the result by applying Stokes' theorem.
\end{snippetexercise}

\begin{snippetsolution}{surfaces-and-flux-exercises-ex15-sol}{Stokes theorem on a circle}
    We first compute the line integral directly. The circle is parametrized counterclockwise by
    \[
        \gamma(t)=(2\cos t,2\sin t,0),
        \qquad 0\leq t\leq 2\pi.
    \]
    Then
    \[
        \gamma'(t)=(-2\sin t,2\cos t,0).
    \]
    Moreover,
    \[
        \mathbf{F}(\gamma(t))
        =
        (4\cos^2 t,2\cos t,2\sin t).
    \]
    Therefore
    \[
        \mathbf{F}(\gamma(t))\cdot \gamma'(t)
        =
        -8\cos^2 t\sin t+4\cos^2 t.
    \]
    Hence
    \[
        \int_\gamma \mathbf{F}\cdot d\mathbf{r}
        =
        \int_0^{2\pi}
        \left(
            -8\cos^2 t\sin t+4\cos^2 t
        \right)
        \,dt.
    \]
    Since
    \[
        \int_0^{2\pi}-8\cos^2 t\sin t\,dt=0
    \]
    and
    \[
        \int_0^{2\pi}4\cos^2 t\,dt=4\pi,
    \]
    we obtain
    \[
        \int_\gamma \mathbf{F}\cdot d\mathbf{r}
        =
        4\pi.
    \]

    We now verify the result using Stokes' theorem. We compute
    \[
        \nabla\times\mathbf{F}
        =
        \begin{vmatrix}
            \mathbf{i} & \mathbf{j} & \mathbf{k} \\
            \partial_x & \partial_y & \partial_z \\
            x^2 & x & y
        \end{vmatrix}
        =
        (1,0,1).
    \]
    Let \(D\) be the disk
    \[
        D=\{(x,y,0):x^2+y^2\leq 4\}.
    \]
    Since the boundary is traversed counterclockwise, we choose the upward unit normal
    \[
        \mathbf{N}=(0,0,1).
    \]
    By Stokes' theorem,
    \[
        \int_\gamma \mathbf{F}\cdot d\mathbf{r}
        =
        \iint_D
        (\nabla\times\mathbf{F})\cdot\mathbf{N}\,dS.
    \]
    Since
    \[
        (\nabla\times\mathbf{F})\cdot\mathbf{N}
        =
        (1,0,1)\cdot(0,0,1)
        =
        1,
    \]
    it follows that
    \[
        \int_\gamma \mathbf{F}\cdot d\mathbf{r}
        =
        \iint_D 1\,dS
        =
        |D|
        =
        4\pi.
    \]
\end{snippetsolution}

\begin{snippetexercise}{surfaces-and-flux-exercises-ex16}{Line integral on an intersection curve}
    Compute \[\int_{\gamma}\left(y^2\,\mathrm{d}x+xy\,\mathrm{d}y+xz\,\mathrm{d}z\right),\] where \(\gamma\) is the intersection curve of the plane \(z=y\) with the cylindrical surface of equation \(x^2+y^2=2x\). Does the result depend on the choice of orientation of \(\gamma\)?
\end{snippetexercise}

\begin{snippetsolution}{surfaces-and-flux-exercises-ex16-sol}{Line integral on an intersection curve}
    The cylindrical surface has equation
    \[
        x^2+y^2=2x,
    \]
    that is
    \[
        (x-1)^2+y^2=1.
    \]
    Therefore the intersection curve can be parametrized by
    \[
        \gamma(t)=(1+\cos t,\sin t,\sin t),
        \qquad 0\leq t\leq 2\pi.
    \]
    Indeed, the condition \(z=y\) is satisfied, and
    \[
        (1+\cos t)^2+\sin^2t
        =
        2+2\cos t
        =
        2(1+\cos t).
    \]

    Along \(\gamma\), we have
    \[
        x=1+\cos t,
        \qquad
        y=\sin t,
        \qquad
        z=\sin t,
    \]
    and
    \[
        dx=-\sin t\,dt,
        \qquad
        dy=\cos t\,dt,
        \qquad
        dz=\cos t\,dt.
    \]
    Hence
    \[
        y^2\,dx+xy\,dy+xz\,dz
        =
        \sin^2t(-\sin t)\,dt
        +(1+\cos t)\sin t\cos t\,dt
        +(1+\cos t)\sin t\cos t\,dt.
    \]
    Therefore
    \[
        y^2\,dx+xy\,dy+xz\,dz
        =
        \left(
            -\sin^3t
            +2\sin t\cos t
            +2\sin t\cos^2t
        \right)\,dt.
    \]
    It follows that
    \[
        \int_\gamma
        \left(
            y^2\,dx+xy\,dy+xz\,dz
        \right)
        =
        \int_0^{2\pi}
        \left(
            -\sin^3t
            +2\sin t\cos t
            +2\sin t\cos^2t
        \right)
        \,dt.
    \]
    Each term has integral equal to zero over \([0,2\pi]\), hence
    \[
        \int_\gamma
        \left(
            y^2\,dx+xy\,dy+xz\,dz
        \right)
        =
        0.
    \]

    Equivalently, since on the plane \(z=y\) we have \(dz=dy\), the form becomes
    \[
        y^2\,dx+xy\,dy+xz\,dz
        =
        y^2\,dx+2xy\,dy
        =
        d(xy^2).
    \]
    Thus the integral over the closed curve \(\gamma\) is zero.

    The result does not depend on the choice of orientation. Reversing the orientation changes the sign of the integral, but since the value is \(0\), it remains \(0\).
\end{snippetsolution}

\begin{snippetexercise}{surfaces-and-flux-exercises-ex17}{Flux through an ellipsoid}
    Compute the outward flux of the field \[\mathbf{F}(x,y,z)=\left(x^3,4y^3,z\right)\] through the surface with equation \[x^2+4y^2+2z^2=4.\]
\end{snippetexercise}

\begin{snippetsolution}{surfaces-and-flux-exercises-ex17-sol}{Flux through an ellipsoid}
    Let
    \[
        E=\{(x,y,z):x^2+4y^2+2z^2\leq 4\}.
    \]
    By the divergence theorem, the outward flux through \(\partial E\) is
    \[
        \iint_{\partial E}\mathbf{F}\cdot \mathbf{N}\,dS
        =
        \iiint_E \operatorname{div}\mathbf{F}\,dV.
    \]
    Since
    \[
        \mathbf{F}(x,y,z)=(x^3,4y^3,z),
    \]
    we have
    \[
        \operatorname{div}\mathbf{F}
        =
        3x^2+12y^2+1.
    \]
    Therefore
    \[
        \Phi
        =
        \iiint_E (3x^2+12y^2+1)\,dV.
    \]

    We transform \(E\) into the unit ball by setting
    \[
        x=2u,\qquad y=v,\qquad z=\sqrt{2}w.
    \]
    Then
    \[
        u^2+v^2+w^2\leq 1,
    \]
    and the Jacobian is
    \[
        \left|\frac{\partial(x,y,z)}{\partial(u,v,w)}\right|
        =
        2\sqrt{2}.
    \]
    Hence
    \[
        \Phi
        =
        2\sqrt{2}
        \iiint_{B(0,1)}
        \left(12u^2+12v^2+1\right)
        \,du\,dv\,dw.
    \]
    By symmetry on the unit ball,
    \[
        \iiint_{B(0,1)}u^2\,dV
        =
        \iiint_{B(0,1)}v^2\,dV
        =
        \frac{4\pi}{15},
    \]
    and
    \[
        \operatorname{Vol}(B(0,1))=\frac{4\pi}{3}.
    \]
    Thus
    \[
        \Phi
        =
        2\sqrt{2}
        \left(
            12\frac{4\pi}{15}
            +
            12\frac{4\pi}{15}
            +
            \frac{4\pi}{3}
        \right).
    \]
    Therefore
    \[
        \Phi
        =
        2\sqrt{2}
        \left(
            \frac{32\pi}{5}
            +
            \frac{4\pi}{3}
        \right)
        =
        2\sqrt{2}\frac{116\pi}{15}
        =
        \frac{232\sqrt{2}\pi}{15}.
    \]
\end{snippetsolution}

\begin{snippetexercise}{surfaces-and-flux-exercises-ex18}{Flux through the boundary of a solid}
    Compute the outward flux of the field \[\mathbf{F}(x,y,z)=\left(z,x^2y,y^2z\right)\] through the surface of the solid \[S=\left\{(x,y,z):x^2+y^2\leq 4,\ 4\sqrt{x^2+y^2}\leq z\leq 1+x^2+y^2\right\}.\]
\end{snippetexercise}

\begin{snippetsolution}{surfaces-and-flux-exercises-ex18-sol}{Flux through the boundary of a solid}
    By the divergence theorem, the outward flux through \(\partial S\) is
    \[
        \iint_{\partial S}\mathbf{F}\cdot \mathbf{N}\,dS
        =
        \iiint_S \operatorname{div}\mathbf{F}\,dV.
    \]
    Since
    \[
        \mathbf{F}(x,y,z)=\left(z,x^2y,y^2z\right),
    \]
    we have
    \[
        \operatorname{div}\mathbf{F}
        =
        \frac{\partial z}{\partial x}
        +
        \frac{\partial x^2y}{\partial y}
        +
        \frac{\partial y^2z}{\partial z}
        =
        x^2+y^2.
    \]

    We use cylindrical coordinates
    \[
        x=r\cos\theta,
        \qquad
        y=r\sin\theta.
    \]
    The solid is described by
    \[
        4r\leq z\leq 1+r^2.
    \]
    This requires
    \[
        4r\leq 1+r^2,
    \]
    that is
    \[
        r^2-4r+1\geq 0.
    \]
    Since also \(0\leq r\leq 2\), we get
    \[
        0\leq r\leq 2-\sqrt{3}.
    \]
    Therefore
    \[
        S=
        \left\{
            (r,\theta,z):
            0\leq r\leq 2-\sqrt{3},
            \ 0\leq\theta\leq 2\pi,
            \ 4r\leq z\leq 1+r^2
        \right\}.
    \]

    Since
    \[
        x^2+y^2=r^2,
        \qquad
        dV=r\,dz\,dr\,d\theta,
    \]
    the flux is
    \[
        \Phi
        =
        \int_0^{2\pi}
        \int_0^{2-\sqrt{3}}
        \int_{4r}^{1+r^2}
        r^2r\,dz\,dr\,d\theta.
    \]
    Hence
    \[
        \Phi
        =
        2\pi
        \int_0^{2-\sqrt{3}}
        r^3(1+r^2-4r)\,dr.
    \]
    Thus
    \[
        \Phi
        =
        2\pi
        \left[
            \frac{r^4}{4}
            +
            \frac{r^6}{6}
            -
            \frac{4r^5}{5}
        \right]_0^{2-\sqrt{3}}.
    \]
    Therefore
    \[
        \Phi
        =
        \frac{\pi}{30}
        \left(
            1392\sqrt{3}-2411
        \right).
    \]
\end{snippetsolution}

\end{document}
